\sectionkicker{Source-backed technical notes}
\subsection*{Paper Watch — agent評価・測定bias・unified multimodal: Technical Notes}
本欄は記事本文で圧縮した一次資料上の情報を、比較・再検証しやすい形で整理したものである。時系列、確認済みの主張、留意点、一次資料URLを掲載する。
\medskip
\subsection*{Theme at a glance}
\addcontentsline{toc}{subsection}{Theme at a glance}
\begin{center}
\footnotesize
\begin{tabularx}{\linewidth}{@{}p{0.20\linewidth}p{0.10\linewidth}p{0.14\linewidth}X@{}}
\toprule
資料 & 位置づけ & 種別 & 時系列 \\
\midrule
AgenticVBench: Can AI Agents Complete Real-World Post-Production Tasks? & 主要資料 & 論文 & 2026-05-26 (論文公開) \\
Identifying and Mitigating Systemic Measurement Bias in Production LLM Inference Benchmarks & 主要資料 & 論文 & 2026-05-22 (論文公開) \\
Representation Forcing for Bottleneck-Free Unified Multimodal Models & 主要資料 & 論文 & 2026-05-29 (論文公開) \\
\bottomrule
\end{tabularx}
\end{center}
\normalsize

\begin{technicalnote}{AgenticVBench: Can AI Agents Complete Real-World Post-Production Tasks?}{主要資料}
\begingroup
% reader-facing Technical Notes break policy
\widowpenalty=10000
\clubpenalty=10000
\displaywidowpenalty=10000
\begin{tabularx}{\linewidth}{@{}>{\bfseries}p{0.22\linewidth}X@{}}
組織 & - \\
種別 & 論文 \\
時系列 & 2026-05-26 (論文公開) \\
\end{tabularx}
\smallskip
{\bfseries 一次資料から整理したtechnical points}
\begin{itemize}[leftmargin=1.5em,itemsep=0.35em]
\item \textbf{一次情報で確認できる事実}: 保存済み一次資料では「AgenticVBench: Can AI Agents Complete Real-World Post-Production Tasks?」が2026-05-26に確認できる。
\end{itemize}
\begin{minipage}{\linewidth}
% reader-facing Technical Notes coherent tail group
\begin{itemize}[leftmargin=1.5em,itemsep=0.35em]
\item \textbf{Author claim}: AgenticVBenchは、text・image・audio・video・planning・tool useを必要とするreal-world post-production workflowでmultimodal agentを評価するbenchmarkを提示している。
\end{itemize}
{\bfseries 読む際の境界}
\begin{itemize}[leftmargin=1.5em,itemsep=0.35em]
\item \textbf{分析上の留意点}: arXivの原メタデータ／abstractは一次論文資料として扱うが、実験結果と一般化可能性はAuthor claimであり、本号では独立再現していない。
\end{itemize}
\begin{samepage}
{\bfseries 一次資料}
\begingroup\sloppy
\begin{itemize}[leftmargin=1.5em,itemsep=0.25em]
\item {\scriptsize\url{http://arxiv.org/abs/2605.27705v1}}
\end{itemize}
\endgroup
\end{samepage}
\end{minipage}
\endgroup
\end{technicalnote}
\medskip

\begin{technicalnote}{Identifying and Mitigating Systemic Measurement Bias in Production LLM Inference Benchmarks}{主要資料}
\begingroup
% reader-facing Technical Notes break policy
\widowpenalty=10000
\clubpenalty=10000
\displaywidowpenalty=10000
\begin{tabularx}{\linewidth}{@{}>{\bfseries}p{0.22\linewidth}X@{}}
組織 & - \\
種別 & 論文 \\
時系列 & 2026-05-22 (論文公開) \\
\end{tabularx}
\smallskip
{\bfseries 一次資料から整理したtechnical points}
\begin{itemize}[leftmargin=1.5em,itemsep=0.35em]
\item \textbf{一次情報で確認できる事実}: 保存済み一次資料では「Identifying and Mitigating Systemic Measurement Bias in Production LLM Inference Benchmarks」が2026-05-22に確認できる。
\end{itemize}
\begin{minipage}{\linewidth}
% reader-facing Technical Notes coherent tail group
\begin{itemize}[leftmargin=1.5em,itemsep=0.35em]
\item \textbf{Author claim}: 論文では、production LLM inference benchmarkにclient-sideのsystemic measurement biasが入り得るとし、Python GILがTTFT/TPOTへ与える影響をqueueing analysisでモデル化している。
\end{itemize}
{\bfseries 読む際の境界}
\begin{itemize}[leftmargin=1.5em,itemsep=0.35em]
\item \textbf{分析上の留意点}: arXivの原メタデータ／abstractは一次論文資料として扱うが、実験結果と一般化可能性はAuthor claimであり、本号では独立再現していない。
\end{itemize}
\begin{samepage}
{\bfseries 一次資料}
\begingroup\sloppy
\begin{itemize}[leftmargin=1.5em,itemsep=0.25em]
\item {\scriptsize\url{http://arxiv.org/abs/2605.24217v2}}
\end{itemize}
\endgroup
\end{samepage}
\end{minipage}
\endgroup
\end{technicalnote}
\medskip

\begin{technicalnote}{Representation Forcing for Bottleneck-Free Unified Multimodal Models}{主要資料}
\begingroup
% reader-facing Technical Notes break policy
\widowpenalty=10000
\clubpenalty=10000
\displaywidowpenalty=10000
\begin{tabularx}{\linewidth}{@{}>{\bfseries}p{0.22\linewidth}X@{}}
組織 & - \\
種別 & 論文 \\
時系列 & 2026-05-29 (論文公開) \\
\end{tabularx}
\smallskip
{\bfseries 一次資料から整理したtechnical points}
\begin{itemize}[leftmargin=1.5em,itemsep=0.35em]
\item \textbf{一次情報で確認できる事実}: 保存済み一次資料では「Representation Forcing for Bottleneck-Free Unified Multimodal Models」が2026-05-29に確認できる。
\item \textbf{Author claim}: Representation Forcingは、unified multimodal model内でvisual representationをautoregressiveなintermediate targetとして明示し、別途pretrained generative VAEを必要としない方式を提案している。
\end{itemize}
{\bfseries 読む際の境界}
\begin{itemize}[leftmargin=1.5em,itemsep=0.35em]
\item \textbf{分析上の留意点}: arXivの原メタデータ／abstractは一次論文資料として扱うが、実験結果と一般化可能性はAuthor claimであり、本号では独立再現していない。
\end{itemize}
\begin{samepage}
{\bfseries 一次資料}
\begingroup\sloppy
\begin{itemize}[leftmargin=1.5em,itemsep=0.25em]
\item {\scriptsize\url{http://arxiv.org/abs/2605.31604v4}}
\end{itemize}
\endgroup
\end{samepage}
\endgroup
\end{technicalnote}
\medskip
